\section{Internalizing the proof in intuitionistic PRA}
\label{sec:internalization}

This section turns the preceding mathematics into a precise arithmetization
plan.  The goal is not merely to say that the proof is ``finitary.''  We
identify the primitive-recursive functions, the certificate relations, the
induction measures, and the final iPRA sentences.

\subsection{The base syntax}

Fix natural-number codes for the following datatypes.

\begin{align*}
 \code{ZCode} & &&\text{signed integers},\\
 \code{QCode} & &&\text{rational numbers},\\
 \code{MonomCode}(v) & &&\text{monomials in $v$ variables},\\
 \code{PolyCode}(v) & &&\text{normalized sparse polynomials},\\
 \code{RatFunCode}(v) & &&\text{fractions of polynomial codes},\\
 \code{PermCode}(d) & &&\text{permutations of $d$ symbols},\\
 \code{GroupCode} & &&\text{finite group tables},\\
 \code{ExprCode} & &&\text{field and radical expressions},\\
 \code{ProofCode} & &&\text{finite derivations},\\
 \code{TowerCode} & &&\text{radical-tower certificates}.
\end{align*}

The constructor tags, list projections, and all recursors are primitive
recursive.  Parsing is total: malformed codes are interpreted as a fixed
error object, so every checker is a total natural-number function.

\subsection{Polynomial and rational-function normalization}

Define the following primitive-recursive functions.

\begin{align*}
 \operatorname{pnorm}&:\code{PolyCode}(v)\to\code{PolyCode}(v),\\
 \operatorname{padd},\operatorname{pmul},\operatorname{pneg}
  &:\code{PolyCode}(v)^2\to\code{PolyCode}(v),\\
 \operatorname{psubst}&:\code{PolyCode}(v)\times
   \code{PolyCode}(w)^v\to\code{PolyCode}(w),\\
 \operatorname{rfeq}&:\code{RatFunCode}(v)^2\to\{0,1\}.
\end{align*}

The sparse normalizer sorts monomials by a primitive-recursive lexicographic
rank, combines adjacent equal exponent vectors, and deletes zero coefficients.
Correctness is proved by quantifier-free induction on the source list.  For
fractions $(P,Q)$ and $(R,S)$,
\[
 \operatorname{rfeq}((P,Q),(R,S))=1
 \quad\Longleftrightarrow\quad
 \operatorname{pnorm}(PS-RQ)=0,
\]
subject to checked nonzero denominators.  This gives decidable equality in the
generic rational-function field $L$.

The elementary symmetric functions $e_1,\ldots,e_5$ and the product
$\prod_i(T-x_i)$ are fixed sparse codes.  A one-time normalization certificate
proves equation~\eqref{eq:generic-coefficients}.

\subsection{Low-degree formula checkers}

For each degree $n\le4$, define a total function
\[
 \operatorname{lowCheck}_n(a,w)\in\{0,1\},
\]
where $a$ is a coefficient tuple and $w$ contains the nonzero and radical
witnesses.  It checks:

\begin{itemize}
\item the leading coefficient is nonzero;
\item every radical power equation;
\item every denominator side condition;
\item the coefficient-by-coefficient factorization identity.
\end{itemize}

The root-list constructor
\[
 \operatorname{lowRoots}_n(a,w)
\]
is primitive recursive.  The core iPRA theorems are equations of the form
\[
 \operatorname{lowCheck}_n(a,w)=1
 \to
 \operatorname{factorCheck}
   \bigl(a,\operatorname{lowRoots}_n(a,w)\bigr)=1.
 \tag{I1}\label{eq:ipra-low}
\]
For $n=2$ the proof is a single normalized identity.  For $n=3$ it uses the
three relations in Section~\ref{sec:cubic}.  For $n=4$ it composes the Ferrari
and quadratic identities.  No induction is needed beyond correctness of the
normalizer.

A separate primitive-recursive branch constructor implements the compatible
radical choices:

\begin{itemize}
\item in the cubic, use the supplied coefficient guard $p=0$ or $p\ne0$;
  the $p=0$ branch adjoins two cube roots and deduces $uv=0$ from
  $(uv)^3=0$;
\item in the quartic, use the supplied coefficient guard $q=0$ or $q\ne0$;
  the first branch is biquadratic and the second proves the Ferrari radical
  $s$ nonzero before division;
\item carry the selected equality or nonzero derivation into the denominator
  checker.
\end{itemize}

Because the guards are values of \code{rfeq}, these are intuitionistically
legitimate case distinctions.

\subsection{Finite permutation groups}

Define
\[
 \operatorname{allPerm}(d)
\]
by bounded enumeration of the $d^d$ lists of length $d$ and filtering by the
bijection predicate.  This is inefficient but primitive recursive.  For
$d=5$ the output has length $120$.  Define composition, inverse, sign,
commutator, subgroup closure, and derived subgroup by bounded folds.

Let \code{S5} and \code{A5} be the resulting fixed codes.  The following are
closed equations between numerals and primitive-recursive terms:
\begin{align}
 |\code{S5}|&=120,\\
 |\code{A5}|&=60,\\
 \operatorname{derived}(\code{S5})&=\code{A5},\\
 \operatorname{derived}(\code{A5})&=\code{A5}.
 \tag{I2}\label{eq:ipra-s5}
\end{align}
They can be proved either by evaluation of the full tables or by the much
smaller explicit commutator certificates of Section~\ref{sec:s5-nonsolvable}.
The latter route is preferable for proof size: verify two products, normality
of a generated subgroup, conjugacy of $3$-cycles, and the generation of $A_5$
by $3$-cycles.

Define the Boolean predicate
\[
 \operatorname{solvableB}(G)=1
\]
by iterating the derived-subgroup function $|G|$ times and testing whether the
identity subgroup has been reached.  Equation~\eqref{eq:ipra-s5} proves
\[
 \operatorname{solvableB}(\code{S5})=0.
 \tag{I3}\label{eq:S5-bool-false}
\]

\subsection{Coding radical-tower solvability certificates}

A tower certificate contains:

\begin{enumerate}
\item a finite list of exponents and radical power relations;
\item the finite lists of conjugate radicands used in each normal-closure
  stage;
\item root-of-unity data;
\item a finite automorphism table for each normal stage;
\item restriction maps and exactness witnesses;
\item embeddings of each abelian kernel in a product of cyclic groups;
\item a final restriction map onto the claimed splitting-field group.
\end{enumerate}

The local checker \code{towerCheck} verifies every polynomial relation and
every finite table equation.  The function
\[
 \operatorname{seriesOfTower}(t)
\]
reads the exact sequences and concatenates the derived-series certificates of
the abelian kernels and quotients.  Its correctness theorem has the form
\[
 \operatorname{towerCheck}(t)=1
 \to
 \operatorname{solvableSeriesCheck}
    (\operatorname{seriesOfTower}(t))=1.
 \tag{I4}\label{eq:tower-to-series}
\]
The proof uses quantifier-free induction on the tower length; the invariant is
an equation saying that the certificate generated for the prefix passes its
checker.  No semantic induction over arbitrary fields occurs inside iPRA.

The extension-of-embeddings steps are also proof-carrying.  A certificate lists
the image chosen for each new generator among a bounded finite root list.  The
checker verifies its defining polynomial and every relation to earlier
generators.  Thus no unbounded search or choice operation is trusted.

\subsection{Coding the generic splitting field action}

The arithmetized generic proof needs only rational-function syntax and the
fixed set of $120$ permutations.

For each permutation code $\pi$, define
\[
 \operatorname{rename}_\pi:\code{RatFunCode}(5)\to\code{RatFunCode}(5)
\]
by renaming exponent-vector coordinates.  The fixed polynomial codes for
$e_j$ satisfy
\[
 \operatorname{rfeq}(\operatorname{rename}_\pi(e_j),e_j)=1
\]
for all $\pi$ and $j$, by a bounded check.  Hence every permutation fixes the
coefficient presentation.

Conversely, an automorphism certificate of the splitting presentation lists
the images of $x_1,\ldots,x_5$.  The factorization code for $P$ and the field
cancellation rules prove that each image is one of the five variable codes.
Distinctness is checked by normalizing $x_i-x_j$.  A bounded injection test
therefore recovers a unique permutation.  The two transformations are inverse
by five rational-function equalities.

The resulting theorem is an isomorphism certificate
\[
 \operatorname{galCode}(P)=\code{S5}.
 \tag{I5}\label{eq:gal-code-s5}
\]
No fixed-field quantification is required in this coding.

\subsection{The arithmetized contradiction}

Let \code{generic5} be the fixed code of $P$, and let
\code{RadSplit(generic5,t)} be the formula asserting that $t$ is a checked
radical splitting certificate.  The radical-tower compiler and
\eqref{eq:tower-to-series} give
\[
 \operatorname{RadSplit}(\code{generic5},t)
 \to
 \operatorname{solvableB}(\operatorname{galCode}(P))=1.
\]
Using~\eqref{eq:gal-code-s5}, this implies
$\operatorname{solvableB}(\code{S5})=1$, contradicting
\eqref{eq:S5-bool-false}.  Therefore iPRA proves
\[
 \forall t\;\neg\operatorname{RadSplit}(\code{generic5},t).
 \tag{I6}\label{eq:ipra-no-radsplit}
\]

A primitive-recursive specialization compiler \code{compileGeneric} maps a
purported universal degree-five formula code $c$ and a successful run witness
$w$ to a radical splitting certificate for \code{generic5}.  Hence
\[
 \forall c\;\bigl(
   \operatorname{UniversalRadicalFormula}(5,c)
   \to\bot\bigr).
 \tag{I7}\label{eq:ipra-abel5}
\]
Padding is a primitive-recursive polynomial operation, so iPRA proves
\[
 \forall n\;\bigl(5\le n\to
   \forall c\;\neg\operatorname{UniversalRadicalFormula}(n,c)\bigr).
 \tag{I8}\label{eq:ipra-abel-all}
\]

\subsection{Transfer to the concise native calculus}

The final formulas~\eqref{eq:ipra-low}, \eqref{eq:ipra-abel5}, and
\eqref{eq:ipra-abel-all} are ordinary iPRA sentences.  The ProveIt file
\code{NativeBiInterpretation.lean}
defines translations between ordinary iPRA formulas and native formulas and
proves \code{native_iff_pra}; the round trips
are provably the identity up to the established equivalence.  Therefore a
Lean proof of the ordinary iPRA derivation immediately yields a derivation in
the selected concise calculus.  No new foundational argument about the tiny
presentation is required.

\subsection{No hidden use of stronger arithmetic}

The proof uses large finite bounds, but proof-theoretic strength is not the
same as computational efficiency.  Factorial, bounded exponentiation, sparse
normalization, finite Gaussian elimination, and finite subgroup closure are
all primitive recursive.  Their correctness proofs use induction on explicit
counters with decidable invariants.  The proof does not require:

\begin{itemize}
\item induction for arbitrary formulas;
\item unbounded minimization;
\item transfinite induction;
\item compactness, Zorn's lemma, or a global algebraic closure;
\item excluded middle for arbitrary field-theoretic propositions.
\end{itemize}

The use of splitting fields and normal closures in the human proof is shorthand
for finite witness records.  In the arithmetized proof, the witnesses are
present in the certificate and every property used later is locally checked.
